\section{Conclusões e Trabalhos Futuros}

\begin{frame}{Conclusões}
  \textbf{O Valor da Engenharia de Software:}
  \begin{itemize}
    \item A negligência documental cobra ``juros'' altos, gerando retrabalhos passíveis de mitigação e opacidade sobre o estado do sistema.
    \item Documentação formal possui importância vital não apenas para futuras manutenções, mas para o aprimoramento contínuo do próprio desenvolvimento.
  \end{itemize}
  
  \vspace{1em}
  
  \textbf{A Contribuição Prática:}
  \begin{itemize}
    \item Entrega de uma documentação de requisitos (Casos de Uso, Regras de Negócio e Dicionários) que mapeia e abrange as reais necessidades do SFM.
    \item Entrega de uma análise de completude dos casos de teste do sistema.
    \item Alerta quanto à necessidade da adoção de práticas da engenharia de software.
  \end{itemize}
\end{frame}

\begin{frame}{Recomendações Práticas}
  \textbf{Mudança no Ciclo de Desenvolvimento:}
  \begin{itemize}
    \item \textbf{Critério Restritivo:} A coordenação deve garantir que nenhuma nova funcionalidade seja codificada sem a sua devida especificação documentada.
  \end{itemize}

  \vspace{1.0em}

  \textbf{Manutenção Contínua:}
  \begin{itemize}
    \item É imperativo atualizar o documento à medida que novas demandas são debatidas.
    \item Qualquer alteração no código-fonte deve ser imediatamente refletida nos requisitos para evitar obsolescência.
  \end{itemize}

  \vspace{1.0em}

  \textbf{Automação de Testes:} 
  \begin{itemize}
      \item Transpor os comportamentos mapeados nos Casos de Uso para uma sintaxe de especificação executável (ex: Gherkin), viabilizando testes automatizados.
  \end{itemize}
\end{frame}

\begin{frame}{Trabalhos Futuros}
    \textbf{Expansão do Escopo Documental:} Elicitação e formalização dos Requisitos Não Funcionais do sistema.
    
    \vspace{1.5em}
    
    \textbf{Análise Prospectiva:} Condução de um estudo comparativo futuro para avaliar de forma empírica os impactos reais na produtividade da equipe e na qualidade do software após a adoção destas práticas.
\end{frame}